\chapter{Evaporative Cooling}
\label{chp:evaporative_cooling}
\fancyhead[C]{Chapter \ref{chp:evaporative_cooling}}

Evaporative cooling is a process that utilizes the principle of heat transfer through the evaporation of water to lower the temperature of an environment. This method is particularly effective in hot and dry climates, where the air has a low humidity level, allowing for more efficient evaporation.

Practical applications for evaporative cooling include residential and commercial cooling systems, industrial processes, and agricultural practices. In residential settings, evaporative coolers, also known as swamp coolers, are commonly used to provide a cost-effective and energy-efficient alternative to traditional air conditioning systems.

\begin{figure}[H]
\centering
\begin{tikzpicture}[scale=0.8, every node/.style={scale=0.8}]

%Duct Casing
\draw [black, dashed] (0,0) -- (0.5,0);
\draw [black] (0.5,0) -- (5,0);
\draw [black, dashed] (5,0) -- (5.5,0);
\draw [black, dashed] (0,4) -- (0.5,4);
\draw [black] (0.5,4) -- (5,4);
\draw [black, dashed] (5,4) -- (5.5,4);

%Air Stream Labels
\node [above] at (0.75,2) {A};
\node [above] at (4.75,2) {B};

%Evaporative Cooler & Air Flow Arrows
\fill [gray, pattern= north west lines, opacity=0.4,draw] (2,0) rectangle (3.5,4);
\draw [->] (1.3,2) -- (1.8,2);
\draw [->] (3.7,2) -- (4.2,2);

%Spray Nozzle and Reservoir
\draw [blue] (2,4.2) -- (2.2,4.4) -- (2.4,4.2);
\draw [blue] (2.6,4.2) -- (2.8,4.4) -- (3.0,4.2);
\draw [blue] (3.2,4.2) -- (3.4,4.4) -- (3.6,4.2);

\draw [black] (1.8,-0.5) -- (1.8,-1) -- (3.7,-1) -- (3.7,-0.5);
\draw [blue, pattern= south west lines] (1.8,-0.6) rectangle (3.7,-1);

\end{tikzpicture}
\caption{Typical schematic representation of an evaporative cooler in a duct or AHU section, the component numbering follows the air path from pre-cooler condition through to the off-cooler condition.}
\label{fig:DuctEvapCooler}
\end{figure}

\section{Principles of Evaporative Cooling}
\label{sec:principles_evaporative_cooling}

Evaporative cooling works on the principle of latent heat of vaporisation, as a psychrometric process the air is cooled along a constant wet bulb line (See Section~\ref{sec:evap_psychrometricprocess}).

Warm dry air is passed over a water-saturated medium, such as a pad or mesh. As the air flows through the medium, water evaporates into the air, absorbing heat from the surrounding environment and lowering the air temperature. The effectiveness of this process is influenced by factors such as ambient temperature, humidity levels, and airflow rates.

\section{Psychrometric Process}
\label{sec:evap_psychrometricprocess}

\todo[inline, color=blue!10]{INSERT PSYCHROMETRIC PROCESS DIAGRAM HERE}

\section{Calculations}
\label{sec:evap_calculations}

Calculation of the cooling effect of evaporative coolers can be performed using equation~\ref{eq:evap_cooling_calculation}, whilst rearranging for $\Theta_{B}$ as shown in~\autoref{der:EvapEfficiency}:

\begin{equation}\label{eq:evap_cooling_calculation}
\eta = \frac{\Theta_{A}-\Theta_{B}}{\Theta_{A}-\Theta_{WB}}
\end{equation}

Where:

\begin{itemize}
    \item $\Theta_{A}$ is the inlet temperature to the evaporative cooler
    \item $\Theta_{B}$ is the outlet temperature of the evaporative cooler
    \item $\Theta_{WB}$ is the wet bulb temperature of the air stream which remains constant
    \item $\eta$ is the efficiency of the evaporative cooler \textit{(Typically between 80-90\%)}
\end{itemize}

This can be re-arranged to give the outlet temperature of the evaporative cooler:

\begin{equation}\label{eq:evap_cooling_calculation_rearanged}
\Theta_{B} = \Theta_{A} - \eta(\Theta_{A} - \Theta_{WB})
\end{equation}

\autoref{eq:evap_cooling_calculation_rearanged} has been derived in \autoref{der:EvapEfficiency} and is used to calculate the outlet temperature of an evaporative cooler given the inlet temperature, wet bulb temperature, and efficiency of the cooler.

\begin{derivation}
\label{der:EvapEfficiency}
\onehalfspacing

This is a derivation of the equation shown previously in \autoref{eq:evap_cooling_calculation}.

First we acknowledge that the evaporative cooling is adiabatic and that the wet-bulb temperature remains constant through the evaporation process.

With an infinitely long evaporative cooler (and an infinitely powerful fan to draw the air through it!), the supply air temperature and wet-bulb temperature are equal. Therefore:

\[\Theta_{B} = \Theta_{WB}\]

Another way to describe this would be to use \textit{Wet Bulb Depression}, this is defined as the difference between dry and wet bulb temperatures, i.e.

\[\Theta_{A} - \Theta_{B}\]

As we know 100\% efficiency is improbable in this situation, an efficiency ratio $\eta$ can be added to the wet bulb depression equation to give the air temperature after the evaporative cooler:

\[\text{Temperature Reduction} = \eta (\Theta_{A} - \Theta_{WB})\]

The achieved temperature after the evaporative cooler $\Theta_{B}$ would then be:

\[\Theta_{B} = \Theta_{A} - \text{Temperature Reduction}\]

Substituting the above Temperature Reduction formula:

\[\Theta_{B} = \Theta_{A} - \eta(\Theta_{A} - \Theta_{WB})\]

\textbf{\textit{QED}}
\end{derivation}

\subsubsection{Evaporative Cooler Example}
\label{ssec:EvapCalExample}

Cardiff summer design conditions (0.4\%): \SI{26.2}{\celsius} dry-bulb and \SI{18.2}{\celsius} wet-bulb.

Wet Bulb Depression at design conditions would be:

$\Theta_{A} - \Theta_{B} = 26.2 - 18.2 = 8$

i.e. A 100\% efficient evaporative cooler would produce a \SI{8}{\celsius} temperature drop at Cardiff design conditions.

Using the formula derived in~\autoref{der:EvapEfficiency} and shown in~\autoref{eq:evap_cooling_calculation_rearanged} can provide the supply air temperature possible with a commercially available evaporative cooler assuming an evaporative cooler efficiency of 85\%:

$\Theta_{B} = 26.2 - 0.85(26.2 - 18.2) = \SI{19.4}{\celsius}$

\section{Commercial Applications}
\label{sec:evap_commercial_applications}

The following subsections provide an overview of the different types of evaporative cooling systems used in commercial applications, including stand-alone evaporative coolers and AHU (Air Handling Unit) evaporative cooling systems.

\subsection{Stand-alone Evaporative Coolers}
\label{ssec:standalone_evap_coolers}

Sometimes nicknamed \textit{'swamp coolers'} or known as direct evaporative cooling, stand-alone evaporative coolers are self-contained units that can be placed in a room or space to provide cooling. These units typically consist of a fan, water reservoir, and a water-saturated medium through which air is drawn. As the air passes through the medium, it is cooled by the evaporation of water. These can be effective in dry climates, but in a sealed room will quickly increase the relative humidity to uncomfortable levels. Their effectiveness is limited in humid climates, where the air already contains a significant amount of moisture; however, in the correct conditions they can provide a cost-effective and energy-efficient cooling solution.

\subsection{AHU Evaporative Cooling}
\label{ssec:ahu_evap_cooling}

Evaporative units can be installed within an \gls{AHU} to cool the supply air to the space, cooling the supply air is the ultimate aim, but the evaporative cooling process can be implemented in different locations with various benefits and drawbacks. The following subsections provide an overview of the different types of AHU evaporative cooling systems.

\subsubsection{Direct AHU Evaporative Cooling}

The evaporative cooler is installed directly in the supply air stream of the \gls{AHU}, cooling the air before it is distributed to the space. This method is effective in reducing the supply air temperature, but it can also increase the humidity of the air, which may not be desirable in certain applications. Direct AHU evaporative cooling is most effective in dry climates where the increase in humidity can be beneficial.

\todo[inline, color=blue!10]{WORKED EXAMPLE OF DIRECT AHU EVAP COOLING HERE}

\subsubsection{Indirect AHU Evaporative Cooling}

This is a cooling only system where outside air is not provided to the space by the AHU. Outside air is passed through an evaporative cooler then over a heat exchanger, and is exhausted. Indoor return air is passed over the heat exchanger, is cooled, then returned to the space.

This is not likely to be a practical method of comfort cooling within the United Kingdom given the relatively humid weather and the need to provide outdoor air to spaces to allow occupants to occupy them safely. There will be specialist use cases for this type of cooler such as data centres, storage facilities, etc where high volumes of outdoor air are not required; however, these types of facilities typically utilise closed loop refrigerant based cooling systems with specialised indoor cooling units such as \GLS{CRAH} units.

\subsubsection{Exhaust Air Stream Evaporative Cooling}

The evaporative cooler is installed in the return air section of the \gls{AHU} allowing the return air stream to be cooled. The air is then passed through a heat exchanger cooling the supply air without the increase in humidity levels as the evaporated water is exhausted to atmosphere once passed through the heat exchanger.

\textbf{Worked Example}

Outside Air conditions: \SI{26.2}{\celsius} dry-bulb and \SI{18.2}{\celsius} wet-bulb.
Room conditions: \SI{22}{\celsius} dry-bulb and 60\% \GLS{RH}

The room conditions indicate a wet-bulb temperature of \SI{15.5}{\celsius} which will be used later to calculate the evaporative cooling effect.

Using~\autoref{eq:evap_cooling_calculation_rearanged}, the return air is passed through an 85\% efficient evaporative cooler and its temperature is reduced to \SI{16.48}{\celsius} as shown in~\autoref{eq:EEAHU_Example}. The \GLS{RH} of the return air stream is increased to 95\%.

\begin{equation}\label{eq:EEAHU_Example}
\Theta_{B} = 22 - 0.85(22 - 15.5) = 16.475
\end{equation}

This cooled return air is then passed through the \gls{AHU} heat exchanger, cooling the supply air to the space without increasing the humidity levels. The supply air temperature can be calculated using the heat exchanger effectiveness and the temperature difference between the cooled return air and the outside air as described in~\autoref{ssec:Vent_HeatRecoveryCalculations}.

The corresponding supply air temperature after the heat exchanger (Assumed 73\% efficient) would be \SI{19.1}{\celsius}, as shown in~\autoref{eq:EEAHU_PHX_Example}. Using the psychrometric chart, the supply air to the space would be \SI{19.1}{\celsius} dry-bulb and \SI{15.8}{\celsius} wet-bulb, which is a reduced humidity level when compared to direct AHU evaporative cooling and is approximately 80-85\% \GLS{RH}.

\begin{equation}\label{eq:EEAHU_PHX_Example}
t_2 = (0.73 \times (16.475 - 26.2)) + 26.2 = 19.1\si{\celsius}, \text{ say } \SI{19}{\celsius}
\end{equation}